\documentclass[a4paper]{article}
% Español (México) con XeLaTeX: fontspec para fuentes Unicode, polyglossia
% para la división silábica y los nombres del documento en la variante mexicana.
\usepackage{fontspec}
\usepackage{polyglossia}
\setdefaultlanguage[variant=mexican]{spanish}
% Los nombres chinos van como «pinyin (original)»: xeCJK da a los caracteres
% Han su propia fuente; sin ella XeLaTeX los omite con un aviso.
% FandolSong no cubre algunos caracteres de nombres (U+73FA); WenQuanYi Zen Hei sí.
\usepackage[AutoFallBack=true]{xeCJK}
\setCJKmainfont{FandolSong-Regular.otf}
\setCJKfallbackfamilyfont{\CJKrmdefault}{WenQuanYi Zen Hei}
% polyglossia no traduce los nombres de listings.
\AtBeginDocument{\providecommand{\lstlistingname}{}\renewcommand{\lstlistingname}{Listado}%
  \providecommand{\lstlistlistingname}{}\renewcommand{\lstlistlistingname}{Índice de listados}}
\usepackage{amsmath, amssymb}
\usepackage{graphicx}
\usepackage[margin=2.35cm]{geometry}
\usepackage[most]{tcolorbox}
\usepackage{booktabs}
\usepackage{tabularx}
\usepackage{longtable}
\usepackage{array}
\usepackage{float}
\usepackage{xcolor}
\usepackage{hyperref}
\usepackage{fancyhdr}

\hypersetup{colorlinks=true, linkcolor=blue!55!black, urlcolor=blue!60!black}
\graphicspath{{./}{figures/}}
\setlength{\parskip}{0.55em}
\setlength{\parindent}{2em}
\setlength{\emergencystretch}{3em}
\renewcommand{\arraystretch}{1.18}

\newtcolorbox{knowledgebox}[1]{enhanced, colback=blue!5!white, colframe=blue!70!black, colbacktitle=blue!70!black, coltitle=white, fonttitle=\bfseries, title={#1}, sharp corners, breakable}
\newtcolorbox{importantbox}[1]{enhanced, colback=yellow!10!white, colframe=yellow!75!black, colbacktitle=yellow!75!black, coltitle=black, fonttitle=\bfseries, title={#1}, sharp corners, breakable}
\newtcolorbox{warningbox}[1]{enhanced, colback=red!5!white, colframe=red!70!black, colbacktitle=red!70!black, coltitle=white, fonttitle=\bfseries, title={#1}, sharp corners, breakable}

\pagestyle{fancy}
\fancyhf{}
\fancyfoot[C]{\thepage}
\fancyfoot[R]{\footnotesize\color{gray}\href{https://github.com/hqhq1025/ai-course-notes}{AI Course Notes}}
\renewcommand{\headrulewidth}{0pt}
\renewcommand{\footrulewidth}{0pt}

\newcommand{\videourl}{https://www.youtube.com/watch?v=bv8ghyTFF9w}
\newcommand{\repourl}{https://github.com/hqhq1025/ai-course-notes}

\begin{document}
\begin{titlepage}
\centering
{\Large Notas de entrevista técnica\par}
\vspace{1.0cm}
{\huge\bfseries Entrevista con Hong Letong (洪乐潼): AI for Math, Lean e intuición matemática\par}
\vspace{0.6cm}
{\Large Zhang Xiaojun conversa con Hong Letong (Carina Hong)\par}
\vspace{0.4cm}
{\large Elaborado a partir del video público de Zhang Xiaojun Podcast, una transcripción generada con GPU y diagramas conceptuales propios\par}
\vspace{0.3cm}
{\large 2026-04-20\par}
\vspace{0.9cm}
\includegraphics[width=0.88\textwidth,height=0.42\textheight,keepaspectratio]{cover.jpg}\par
\vfill
\begin{tcolorbox}[width=0.92\textwidth, colback=black!2!white, colframe=black!55, sharp corners]
\textbf{Título del video}: 137. Entrevista de 4 horas con Hong Letong: AI for Math, convertir las matemáticas a Lean, las demostraciones de «El Libro» de las matemáticas, la intuición, lo creado y lo descubierto\par
\textbf{Canal del video}: Zhang Xiaojun Podcast\par
\textbf{Duración del video}: 04:24:17\par
\textbf{Enlace del video}: \href{\videourl}{\nolinkurl{\videourl}}\par
\textbf{Descripción del material}: YouTube no ofrece subtítulos; el cuerpo del texto se elaboró a partir de la transcripción con faster-whisper large-v3 en CUDA, los capítulos del video, la portada y figuras didácticas propias. Las tomas fijas de la entrevista no se reutilizan en el cuerpo del texto.\par
\textbf{Página del proyecto}: \href{\repourl}{\nolinkurl{\repourl}}\par
\end{tcolorbox}
\end{titlepage}

\tableofcontents
\newpage
\section{Introducción: por qué este episodio es la clase de entrada a AI for Math}

La protagonista de esta entrevista es Hong Letong (洪乐潼), Carina Hong, fundadora de Axiom Math, cuya línea de trabajo es AI for Math: que la IA participe en la resolución de problemas matemáticos, la formalización de demostraciones, la búsqueda de demostraciones y la colaboración en la investigación matemática. El título del video contiene muchas palabras clave: Lean, las demostraciones de «El Libro» de las matemáticas, intuición, ser creado o ser descubierto, Axiom, Ken Ono, Neo Lab. Todas ellas apuntan a una misma pregunta: si la IA ha de entrar en la investigación matemática real, no basta con que aprenda a resolver problemas; también tiene que comprender el lenguaje matemático, la estructura de las demostraciones, los sistemas formales y la intuición de los matemáticos.

\begin{importantbox}{Tesis central del episodio}
Lo decisivo en AI for Math no es que el modelo resuelva más problemas de competencia, sino cerrar el ciclo que une la intuición matemática, las demostraciones informales, el lenguaje formal, el proof assistant y la verificación por máquina. Lean no es un adorno: es la interfaz fundamental que convierte lo que «parece correcto» en algo «verificable por máquina».
\end{importantbox}

\begin{knowledgebox}{Nota sobre la estrategia visual}
Este video es una entrevista con plano fijo; no tiene slides didácticas, pizarrón ni demostraciones de producto. Conforme al estándar de pódcast de este repositorio, el cuerpo del texto no vuelve a insertar fotogramas de las personas; la portada sirve para identificar la fuente, y el contenido didáctico del cuerpo se apoya en diagramas conceptuales, mapas de ruta y glosarios.
\end{knowledgebox}

\subsection{Resumen de la sección}

Este episodio no es una entrevista común sobre emprendimiento, sino una entrada conceptual a AI for Math: aborda a la vez la filosofía de las matemáticas, las demostraciones formales, Lean, la ruta de emprendimiento, la intuición en la investigación y la estructura de recursos de Neo Lab.
\section{Guía introductoria: cómo juzgar si un sistema de AI for Math es confiable}

Al estudiar este episodio, la pregunta «¿es confiable un sistema de AI for Math?» puede dividirse en seis preguntas. Primera: ¿puede comprender con precisión los objetos matemáticos y los enunciados de los teoremas, en lugar de limitarse a hacer una paraphrase en lenguaje natural? Segunda: ¿puede convertir una demostración informal en una forma que un proof assistant como Lean/Coq pueda verificar? Tercera: cuando una demostración falla, ¿puede leer el goal state y repararla, en lugar de reescribir un texto alucinado? Cuarta: ¿puede aprovechar una biblioteca de teoremas, en lugar de empezar desde cero cada vez? Quinta: ¿puede acelerar la investigación de los matemáticos, y no solo resolver problemas ya conocidos? Sexta: ¿puede explicar por qué una demostración es significativa?

\begin{importantbox}{El método de las seis preguntas}
Si un sistema de AI for Math solo cumple las dos primeras preguntas, se parece más a una herramienta de formalization; si cumple la tercera y la cuarta, empieza a tener capacidades de Agent; solo si cumple la quinta y la sexta se acerca a ser un colaborador de investigación matemática.
\end{importantbox}

\subsection{Por qué importan estas seis preguntas}

Convierten la pregunta difusa «¿sabe matemáticas la IA?» en capacidades verificables. Comprender los objetos corresponde a la capacidad semántica; la formalización corresponde a la capacidad de lenguaje y de sistema de tipos; leer el goal state corresponde a la retroalimentación del entorno; aprovechar la biblioteca de teoremas corresponde a la recuperación y la memoria; acelerar la investigación corresponde al valor de producto; explicar el significado corresponde a la sensibilidad estética matemática. Un AI for Math de alto nivel requiere que todas estas capacidades se presenten juntas.

\subsection{Resumen de la sección}

Estas seis preguntas atraviesan el resto de la nota. Aun sin conocer de antemano los detalles de Lean, el lector puede usarlas para juzgar si las distintas líneas de AI for Math están resolviendo problemas, construyendo demostraciones o intentando entrar en la investigación propiamente dicha.
\section{¿Las matemáticas se descubren o se crean?}

Al inicio de la entrevista se discute «lo creado y lo descubierto». Este es el trasfondo filosófico de AI for Math. Si las matemáticas se descubren, la IA necesita aprender estructuras objetivas: las relaciones estables entre números primos, simetrías, topología, curvatura y demostraciones. Si las matemáticas se crean, la IA necesita aprender cómo los seres humanos eligen definiciones, axiomas, notaciones y lenguajes teóricos. En la realidad ambas cosas se entrelazan: los objetos matemáticos tienen restricciones objetivas, mientras que el lenguaje matemático y los caminos de demostración los crean las personas.

\begin{figure}[H]
\centering
\includegraphics[width=0.92\textwidth]{math-created-discovered.png}
\caption{Dos perspectivas de las matemáticas: la estructura descubierta y el lenguaje creado.}
\end{figure}

\begin{knowledgebox}{Lectura de la figura: por qué AI for Math enfrenta dos niveles de tarea a la vez}
El lado izquierdo destaca la objetividad de la estructura matemática; el lado derecho, el carácter creativo del lenguaje matemático y de los sistemas axiomáticos. AI for Math no puede limitarse a la coincidencia de patrones ni a generar texto elegante; debe captar estructuras estables y, al mismo tiempo, escribir el razonamiento en un lenguaje verificable.
\end{knowledgebox}

\subsection{bounded attention y free attention}

En la entrevista aparece la expresión bounded vs free attention, que puede entenderse como dos modos del pensamiento matemático: la bounded attention avanza localmente en torno a definiciones, condiciones, objetivos y herramientas existentes; la free attention, en cambio, permite la asociación, la estética, la analogía y los saltos. La investigación matemática necesita alternar entre ambas: la demostración formal se inclina hacia lo bounded, y la generación de conjeturas e intuiciones, hacia lo free.

\begin{figure}[H]
\centering
\includegraphics[width=0.92\textwidth]{bounded-free-attention.png}
\caption{La bounded attention y la free attention en la intuición matemática.}
\end{figure}

\begin{knowledgebox}{Lectura de la figura: a qué sirve cada tipo de atención}
La bounded attention es adecuada para la revisión de demostraciones, el razonamiento local y la formalización en Lean; la free attention es adecuada para descubrir patrones, plantear conjeturas y establecer analogías entre distintos campos. Un sistema sólido de AI for Math necesita que ambas colaboren, en lugar de saber solo hacer búsquedas rigurosas o solo hacer asociaciones divergentes.
\end{knowledgebox}

\begin{warningbox}{No hay que contraponer la intuición y la formalización}
La intuición matemática no es enemiga de la formalización. Un buen sistema de formalización convierte la intuición en objetos verificables; una buena intuición, a su vez, ayuda a elegir los problemas y las rutas de demostración que vale la pena formalizar.
\end{warningbox}

\subsection{Resumen de la sección}

La dificultad de AI for Math proviene de una doble estructura: las matemáticas se parecen tanto a un mundo objetivo como a una obra de ingeniería del lenguaje humano. El modelo debe alternar continuamente entre la asociación libre y la verificación rigurosa.
\section{La pila tecnológica de AI for Math: de la intuición a la demostración verificable}

Muchas personas reducen AI for Math a «si el modelo sabe resolver problemas de matemáticas». Esa es solo la capa más superficial. El objetivo real es que la IA parta de problemas, conjeturas y razonamiento informal, y avance paso a paso hacia enunciados formales, búsqueda de demostraciones, verificación por máquina y retroalimentación humana.

\begin{figure}[H]
\centering
\includegraphics[width=0.94\textwidth]{ai-for-math-stack.png}
\caption{La pila tecnológica de AI for Math: del problema a la verificación y, después, a la retroalimentación de los matemáticos.}
\end{figure}

\begin{knowledgebox}{Lectura de la figura: AI for Math no es una capacidad aislada del modelo}
Cada capa de la figura puede convertirse en un cuello de botella: comprender el problema requiere semántica y conocimiento de contexto; el razonamiento informal requiere intuición; la formalización requiere los lenguajes Lean/Coq; la búsqueda de demostraciones requiere una biblioteca de teoremas y tactics; la verificación requiere un proof assistant; por último, se necesita que los matemáticos juzguen la dirección y el significado del trabajo.
\end{knowledgebox}

\subsection{Términos clave: vocabulario de AI for Math}

\noindent\begin{tabularx}{\textwidth}{p{0.22\textwidth}p{0.32\textwidth}X}
\toprule
Término & Problema que resuelve & Relación con este episodio \\
\midrule
Lean & Lenguaje de matemáticas formales y proof assistant & Convierte demostraciones en lenguaje natural en demostraciones que una máquina puede comprobar. \\
Coq & Otro tipo de proof assistant & Muestra que la demostración formal no tiene a Lean como única vía. \\
Tactic & Paso de demostración automático o semiautomático & Permite que el modelo construya paso a paso demostraciones comprobables. \\
Theorem Proving & Demostración de teoremas & Una de las tareas centrales de AI for Math. \\
Verifier & Verificador & Determina si una demostración es válida y evita la alucinación de algo que «parece correcto». \\
Formalization & Formalización & Lleva artículos, definiciones y conjeturas a un sistema verificable. \\
\bottomrule
\end{tabularx}

\subsection{Por qué el proof assistant es decisivo}

Las demostraciones matemáticas en lenguaje natural suelen omitir pasos y dependen del contexto y la intuición del lector. Un proof assistant exige que cada paso cumpla las reglas. Para la IA, esto constituye una retroalimentación fuerte: la demostración generada o bien la acepta el kernel, o bien falla y devuelve el estado del objetivo. Esta retroalimentación es más rigurosa que una calificación subjetiva y resulta más adecuada para el entrenamiento y la evaluación.

\begin{importantbox}{La verificabilidad es la ventaja defensiva de AI for Math}
Las matemáticas son uno de los pocos campos de alto valor capaces de ofrecer señales de verificación fuertes. Sistemas como Lean/Coq hacen que la salida del modelo no solo «se parezca a una demostración», sino que deba superar una comprobación por máquina. Esto le da a AI for Math, al mismo tiempo, valor científico y capacidad de entrenarse desde la ingeniería.
\end{importantbox}

\subsection{Resumen de la sección}

La pila tecnológica de AI for Math abarca la intuición, los lenguajes formales, la búsqueda de demostraciones y la verificación por máquina. La verdadera dificultad está en conectar estas capas, no en resolver problemas en una sola de ellas.
\section{Convertir las matemáticas a Lean: el ciclo cerrado de la demostración formal}

Convertir las matemáticas a Lean no es una simple traducción. Las demostraciones en lenguaje natural suelen omitir el contexto, saltarse pasos, apoyarse en la intuición geométrica o citar teoremas de forma imprecisa. Lean exige explicitar el tipo de cada objeto, el enunciado del teorema, las dependencias, las tactic y la validez de cada paso. Esto le da a la IA tanto oportunidades como restricciones.

\begin{figure}[H]
\centering
\includegraphics[width=0.92\textwidth]{lean-proof-loop.png}
\caption{Ciclo cerrado desde la demostración en lenguaje natural hasta la verificación en Lean.}
\end{figure}

\begin{knowledgebox}{Lectura de la figura: por qué el ciclo cerrado de Lean es adecuado para el entrenamiento}
La demostración en lenguaje natural se convierte primero en un formal statement; después se generan las tactic o el proof term, y el kernel de Lean comprueba si son válidos. Si fallan, el sistema devuelve el error/goal state, y una persona o el modelo hace la corrección a partir de ello. Este ciclo cerrado proporciona una retroalimentación de entrenamiento más sólida que la generación de texto ordinaria.
\end{knowledgebox}

\subsection{Tres tipos de dificultades de la formalización}

La primera es la dificultad semántica: expresiones del lenguaje natural como «es evidente» o «de aquí se sigue» deben desarrollarse por completo. La segunda es la dificultad de las dependencias de la biblioteca: en Lean, los nombres y la forma de uso de los teoremas, las definiciones y los lemma deben ser exactos. La tercera es la dificultad de búsqueda: el espacio de demostraciones es enorme, y el modelo debe decidir qué subobjetivo demostrar primero, qué lemma invocar y cuándo reescribir el objetivo.

\begin{warningbox}{Lean no hace que el modelo entienda matemáticas por sí solo}
Lean puede verificar si una demostración es válida, pero no elige por el modelo los problemas importantes ni le proporciona por sí solo un sentido estético matemático. AI for Math también necesita la intuición, la selección de problemas y el criterio teórico de los matemáticos humanos.
\end{warningbox}

\subsection{Las demostraciones del Libro celestial de las matemáticas}

El título de la entrevista menciona «las demostraciones del Libro celestial de las matemáticas», que pueden entenderse como aquellas demostraciones elegantes y breves que parecen salidas de un libro celestial. Si la IA solo hace búsqueda por fuerza bruta, puede producir demostraciones largas; solo si aprende la intuición y el sentido estético matemáticos podrá acercarse a este tipo de demostraciones. La verificación formal garantiza la corrección, pero la elegancia sigue requiriendo un juicio estético.

\begin{importantbox}{Corrección y elegancia}
La verificación por máquina resuelve si una demostración «es correcta», pero no resuelve directamente si «es buena». La investigación matemática también se interesa en si la demostración explica el fenómeno, si revela una estructura y si es generalizable. El objetivo a largo plazo de AI for Math debe buscar a la vez la corrección, la legibilidad y la comprensión profunda.
\end{importantbox}

\subsection{Resumen de la sección}

Lean convierte las demostraciones en objetos ejecutables y verificables. Lo esencial de AI for Math es conectar la intuición del lenguaje natural con el ciclo cerrado de verificación por máquina, sin perder el sentido estético matemático.
\section{De la resolución de problemas a la investigación matemática}

AI for Math puede dividirse en distintos niveles: resolver problemas, demostración formal, completado de demostraciones, generación de conjeturas y colaboración en investigación. Cuanto más se avanza, más se acerca a la investigación matemática real y más difícil es evaluar. Los problemas de concurso tienen respuesta, las demostraciones en Lean tienen un verifier, mientras que en los problemas de investigación hay que juzgar «si vale la pena abordarlos».

\begin{figure}[H]
\centering
\includegraphics[width=0.96\textwidth]{ai-for-math-roadmap.png}
\caption{Hoja de ruta de AI for Math: de los problemas matemáticos a la investigación matemática.}
\end{figure}

\begin{knowledgebox}{Lectura de la figura: por qué es más difícil cuanto más a la derecha}
La resolución de problemas suele tener una respuesta; la demostración formal se verifica con un proof assistant; el completado de demostraciones tiene objetivos locales; en cambio, la generación de conjeturas y la colaboración en investigación requieren sentido estético, criterio para elegir direcciones y una estructura teórica de largo plazo. Cuanto más cerca de la investigación, menos se puede depender solo de un benchmark.
\end{knowledgebox}

\subsection{La intuición del matemático}

La segunda mitad de la entrevista trata sobre «la intuición del matemático». La intuición no es algo místico, sino un juicio condensado sobre estructuras, contraejemplos, rutas de demostración y estética, formado tras un entrenamiento prolongado. La IA puede aprender patrones a partir de grandes cantidades de datos y de retroalimentación, pero que llegue a formar algo parecido a la intuición de un matemático depende de si logra integrar la retroalimentación formal, el razonamiento informal y los objetivos de investigación.

\begin{knowledgebox}{La intuición puede entrenarse, pero no equivale a ajustar datos}
La intuición matemática surge de la elección de problemas, la experiencia con los fracasos, las analogías estructurales y el juicio estético. La IA puede aprender de bibliotecas de demostraciones, artículos, retroalimentación interactiva y correcciones de matemáticos, pero necesita un ciclo cerrado de largo plazo, no datos supervisados de una sola vez.
\end{knowledgebox}

\subsection{Un viaje a la Luna o tiene éxito o fracasa}

La expresión de la entrevista «un viaje a la Luna o tiene éxito o fracasa» refleja el carácter de alto riesgo del emprendimiento en AI for Math. La investigación matemática no es como una aplicación común, en la que se pueden validar indicadores comerciales en pasos pequeños: o impulsa de verdad la capacidad de demostrar y descubrir, o es solo una demo vistosa. La dificultad de un Neo Lab como Axiom consiste en que debe ocuparse a la vez de la ciencia, el producto, el equipo y el financiamiento.

\begin{warningbox}{La evaluación de AI for Math no puede basarse solo en el financiamiento}
Las grandes rondas de financiamiento y la incorporación de matemáticos de primer nivel son señales, pero no resultados. Los resultados reales dependen de si se producen demostraciones verificables, se acelera el trabajo de los matemáticos y se forman bibliotecas de teoremas reutilizables y modelos de colaboración en investigación.
\end{warningbox}

\subsection{Resumen de la sección}

La ruta de AI for Math va de la resolución de problemas a la colaboración en investigación. Cuanto más se acerca a la investigación, más necesita la intuición del matemático, sistemas formales y retroalimentación de largo plazo, y no un único benchmark.
\section{Axiom y los Neo Lab: qué requiere emprender en AI for Math}

La descripción del video menciona que Axiom cerró una ronda Serie A de 200 millones de dólares con una valuación de 1,600 millones de dólares; también llamó la atención que Ken Ono dejara o pusiera en pausa la trayectoria tradicional de profesor para unirse a Axiom. Más allá de cómo lo narren los medios, esto muestra que AI for Math pasó de ser un problema académico a ser un problema de Neo Lab con alta densidad de capital.

\begin{figure}[H]
\centering
\includegraphics[width=0.92\textwidth]{ai-math-neolab.png}
\caption{Componentes de un Neo Lab de AI for Math: matemáticos, sistemas de formalización, entrenamiento de modelos, financiamiento y capacidad de cómputo, y objetivos de producto.}
\end{figure}

\begin{knowledgebox}{Lectura de la figura: un Neo Lab no es una empresa emergente común}
Los cinco elementos de la figura apuntan en conjunto a organizaciones como Axiom: la red de matemáticos de primer nivel aporta los problemas y el criterio estético, Lean/Coq aporta la base de formalización, el entrenamiento de modelos aporta la capacidad de automatización, el financiamiento y la capacidad de cómputo sostienen los experimentos a escala, y los objetivos de producto convierten la investigación en sistemas utilizables.
\end{knowledgebox}

\subsection{Por qué importa la red de matemáticos}

Los datos matemáticos no son datos web comunes. Los problemas realmente valiosos, las rutas de demostración, la organización de las bibliotecas de teoremas y el juicio estético suelen estar en manos de la comunidad de matemáticos. Que se sume un matemático como Ken Ono no es solo un respaldo: es probable que aporte la selección de problemas, una red de investigación y criterios de verificación.

\begin{importantbox}{El círculo virtuoso de datos de AI for Math}
Si el sistema puede ayudar a los matemáticos a formalizar demostraciones, completar lemmas y encontrar contraejemplos, y los matemáticos a su vez pueden devolver correcciones y evaluaciones, podría formarse un círculo virtuoso de datos. Sin la participación de matemáticos, el modelo tiende a quedarse en bancos de problemas y demostraciones superficiales.
\end{importantbox}

\subsection{La emprendedora menos probable}

En la entrevista, Hong Letong (洪乐潼) habla de «la emprendedora menos probable». Este tipo de Neo Lab es distinto del emprendimiento tradicional de aplicaciones: quien lo funda debe entender matemáticas, IA, sistemas de formalización, organización de la investigación y financiamiento. Se parece más a comprimir un instituto de investigación, un equipo de modelos y una empresa de producto en una sola organización de alta densidad.

\begin{warningbox}{El doble riesgo del emprendimiento orientado a la investigación}
Si es demasiado académico, la ruta hacia el producto y los ingresos no queda clara; si se orienta demasiado al producto, puede perder profundidad en la investigación básica. AI for Math necesita mantener la tensión entre los objetivos científicos y la implementación de ingeniería.
\end{warningbox}

\subsection{Resumen de la sección}

Axiom representa un nuevo tipo de AI Lab: los problemas de investigación son lo bastante fundamentales, la densidad de capital es lo bastante alta y el equipo necesita conectar a la vez a los matemáticos, los sistemas de formalización y el entrenamiento de modelos. No es un emprendimiento SaaS común ni un laboratorio académico tradicional.
\section{Términos clave: índice de conceptos de este episodio}

\begin{longtable}{p{0.23\textwidth}p{0.32\textwidth}p{0.36\textwidth}}
\toprule
Término & Explicación en una frase & Función en este episodio \\
\midrule
AI for Math & Uso de la IA para apoyar la resolución, la demostración, la formalización y la colaboración en investigación matemática & Dirección central de todo el texto. \\
Lean & Lenguaje de matemáticas formales y proof assistant & Convierte las demostraciones en objetos verificables por máquina. \\
Proof Assistant & Sistema de software que ayuda a escribir y verificar demostraciones formales & Proporciona un verifier riguroso. \\
Formalization & Conversión de las matemáticas en lenguaje natural a un sistema formal & Paso decisivo para conectar la IA con las bibliotecas matemáticas. \\
Tactic & Comando de estrategia que hace avanzar una demostración en un proof assistant & Acción que el modelo puede generar o buscar. \\
Verifier & Mecanismo que comprueba si una salida es válida & Aporta retroalimentación estricta al entrenamiento y a la evaluación. \\
Bounded Attention & Atención local acotada por el objetivo y las definiciones & Sostiene la demostración rigurosa y la formalización. \\
Free Attention & Asociación abierta y juicio estético & Sostiene la formulación de conjeturas y las nuevas rutas de demostración. \\
Neo Lab & Nuevo tipo de laboratorio de IA con alta densidad de capital y orientado a la investigación & Tipo de organización de Axiom. \\
\bottomrule
\end{longtable}

\subsection{Resumen de la sección}

Los términos de este episodio giran en torno a un eje central: cómo llevar las matemáticas de la intuición humana y el lenguaje natural a sistemas de IA verificables, entrenables y que permitan la colaboración.

\section{Conexión con los dos episodios anteriores}

La entrevista con Su Yu (苏煜) en EP139 trata la historia técnica de los Agent; la entrevista con Luo Fuli (罗福莉) en EP138 trata cómo el paradigma Agent transforma los laboratorios de modelos, y la entrevista con Hong Letong (洪乐潼) en EP137 lleva la frontera de la IA hasta la investigación matemática. Leídos en conjunto, los tres episodios forman un mapa continuo que va de los agentes digitales de propósito general y la organización de los laboratorios de modelos a los sistemas de descubrimiento científico.

\begin{importantbox}{El hilo conductor de los tres episodios leídos en conjunto}
EP139 explica por qué los Agent se convirtieron en un nuevo paradigma; EP138 explica cómo los equipos de modelos reorganizan el posentrenamiento y la capacidad de cómputo en función de los Agent; EP137 explica cómo la IA entra en las matemáticas, un campo altamente formalizado y verificable que, sin embargo, depende de la intuición.
\end{importantbox}

\subsection{Resumen de la sección}

AI for Math es una rama científica de alto valor en la era de los Agent. Reúne en un mismo problema el uso de herramientas, la verificación formal, el razonamiento de largo plazo y la colaboración con expertos humanos.
\section{Aguas profundas: intuición matemática, el ecosistema de Lean y la colaboración en la investigación}

Los capítulos anteriores desglosaron AI for Math en una pila tecnológica, pero lo verdaderamente interesante de la entrevista con Hong Letong (洪乐潼) es que ella no concibe las matemáticas como un simple «conjunto de problemas». En la investigación matemática hay estética, hay fracasos, hay dedicación a largo plazo y también hay reconocimiento de la comunidad. Si la IA solo sabe generar respuestas dentro de bancos de problemas ya conocidos, sigue muy lejos de la investigación matemática; solo cuando puede ayudar a los matemáticos a descubrir estructuras, formalizar demostraciones, completar los lemma faltantes y verificar contraejemplos entra de verdad en el proceso de investigación.

\subsection{Por qué las matemáticas no son una tarea común de NLP}

Las tareas comunes de NLP suelen admitir aproximaciones semánticas: un resumen puede redactarse de distintas maneras, una traducción puede tener varias expresiones equivalentes y las respuestas de un chat también admiten diferencias de estilo. Las demostraciones matemáticas son distintas: basta con omitir una condición clave para que toda la demostración pueda fallar. Los sistemas formales llevan este rigor al extremo y exigen que cada objeto, tipo, dependencia y paso de razonamiento sea válido.

\begin{importantbox}{La particularidad de las tareas matemáticas}
Las matemáticas tienen a la vez fuertes restricciones formales y una fuerte creatividad. Las restricciones formales provienen del proof assistant y de la verificación de teoremas; la creatividad proviene de las conjeturas, la estética, la elección de definiciones y el diseño de la ruta de demostración. AI for Math debe atender ambos extremos al mismo tiempo.
\end{importantbox}

\subsection{Por qué es importante el ecosistema de Lean}

El valor de Lean no reside solo en ser un verificador, sino en su ecosistema: la biblioteca de teoremas mathlib, las contribuciones de la comunidad, el estilo de formalización, la tradición de tactic, los materiales didácticos y la participación de los matemáticos. Para escribir demostraciones en Lean, un modelo necesita entender qué hay en la biblioteca, cómo se nombran los teoremas, qué tipos de tactic se usan con frecuencia y cómo interpretar los objetivos fallidos. Esto se parece mucho al autocompletado de código común, pero con restricciones más estrictas.

\begin{knowledgebox}{Semejanzas y diferencias entre Lean y el código}
La semejanza está en que ambos tienen sintaxis, tipos, bibliotecas, mensajes de error y retroalimentación ejecutable; la diferencia está en que el objetivo de Lean no es ejecutar un programa, sino demostrar proposiciones. Un programa puede apoyarse en pruebas que cubran los casos comunes; una demostración, en cambio, debe cumplirse en todos los casos.
\end{knowledgebox}

\subsection{El espacio de búsqueda de Proof Search}

La búsqueda de demostraciones no es una generación lineal. Un teorema puede requerir construir primero un lemma intermedio, elegir la variable de inducción, transformar la forma del objetivo, invocar teoremas existentes o demostrar primero que no existe un contraejemplo. El espacio de búsqueda es enorme y los caminos erróneos son muchos. La ventaja de la IA es que puede hacer intentos en paralelo y aprovechar conocimiento semántico previo; su desventaja es que tiende a generar pasos que parecen razonables pero que no están disponibles en la biblioteca.

\begin{warningbox}{Fallas comunes de la búsqueda de demostraciones}
El modelo puede conocer la idea matemática pero no saber cómo se llama el teorema correspondiente en la biblioteca de Lean; también puede conocer la tactic pero convertir el objetivo en una forma más difícil; o incluso puede generar una demostración que es válida en lenguaje natural pero a la que le faltan premisas en el sistema formal. La propia información de las fallas es, en sí misma, datos de entrenamiento.
\end{warningbox}

\subsection{Resumen de la sección}

Lo esencial de AI for Math no es «un modelo un poco más inteligente», sino descomponer la investigación matemática en un ciclo cerrado con retroalimentación: la intuición propone la dirección, la formalización restringe el proceso, el verificador proporciona retroalimentación estricta y el matemático juzga el significado.
\section{Juicio clave sobre la ruta de Axiom: por qué crear una empresa de investigación}

Un Neo Lab como Axiom no es una empresa de aplicaciones común. Enfrenta una paradoja de producto muy profunda: la investigación matemática tiene muy pocos usuarios finales y su comercialización a corto plazo es menos clara que la de las aplicaciones de IA de propósito general; pero, si el sistema llega a acelerar de verdad la investigación matemática, podría influir en la ciencia, la ingeniería, la criptografía, la física, los materiales y la propia IA. El alto riesgo y el alto potencial coexisten.

\subsection{La señal del caso Ken Ono}

A los medios les gusta contar la historia de «un profesor con plaza definitiva de 57 años que trabaja para un fundador de 24 años», pero lo que merece más atención es la señal en sí: que matemáticos de primer nivel estén dispuestos a dedicar su tiempo a AI for Math indica que no lo consideran un juguete. Lo que aporta la incorporación de matemáticos no es valor publicitario, sino selección de problemas, criterio estético de investigación, juicio sobre la calidad de las demostraciones y vínculos con la comunidad.

\begin{importantbox}{Por qué los matemáticos son insustituibles}
La IA puede generar demostraciones candidatas, pero son los matemáticos quienes deciden qué problemas son importantes, qué demostraciones tienen poder explicativo y qué direcciones merecen inversión. Si AI for Math se desvincula de la comunidad matemática, puede degenerar fácilmente en la automatización de un banco de problemas.
\end{importantbox}

\subsection{El financiamiento no es un resultado, sino runway}

La gran ronda de financiamiento de Axiom muestra que el capital está dispuesto a apostar por AI for Math, pero el financiamiento es solo runway. Lo que realmente debe validarse es si puede construir contribuciones sostenibles a la biblioteca de teoremas, herramientas de automatización de demostraciones, flujos de trabajo para matemáticos y productos de colaboración en investigación. De lo contrario, el monto del financiamiento solo amplificará la presión.

\begin{warningbox}{Los riesgos de un Neo Lab}
Una empresa de investigación corre el riesgo de no encajar en ninguno de los dos lados: la academia la considera demasiado comercial y el mundo empresarial, demasiado básica. Para atravesar ese hueco necesita un producto intermedio verificable, por ejemplo, una mayor eficiencia en las demostraciones formales, el crecimiento de la biblioteca de teoremas, una plataforma de colaboración para matemáticos o un avance de investigación de alto valor.
\end{warningbox}

\subsection{Resumen de la sección}

La ruta de Axiom no consiste en «hacer un robot conversacional de matemáticas», sino en intentar organizar a matemáticos, sistemas formales y entrenamiento de modelos en una máquina de investigación. Tanto su riesgo como su valor provienen de ese posicionamiento.
\section{De las matemáticas al descubrimiento científico en sentido amplio}

AI for Math también es importante porque las matemáticas son la parte más formalizada y verificable del descubrimiento científico. Si la IA logra cerrar en matemáticas el ciclo «formular conjeturas--construir demostraciones--verificación por máquina--evaluación humana», métodos similares podrían trasladarse a la física, los materiales, la verificación de programas y el descubrimiento de algoritmos.

\subsection{Por qué las matemáticas son un campo de pruebas}

Las matemáticas tienen tres ventajas. Primera: la corrección puede verificarse de forma rigurosa mediante sistemas formales. Segunda: los datos históricos y las bibliotecas de teoremas se han acumulado de manera estructurada. Tercera: la comunidad matemática tiene una tradición de evaluación de la elegancia, el poder explicativo y la importancia que es clara, pero difícil de cuantificar. Esto hace que sean adecuadas para el entrenamiento de máquinas y, a la vez, que pongan en evidencia las insuficiencias de ese entrenamiento.

\begin{knowledgebox}{Tres tipos de retroalimentación en el descubrimiento científico}
Las matemáticas aportan verificación formal; el código aporta pruebas y retroalimentación de ejecución; las ciencias experimentales aportan retroalimentación de observaciones y experimentos. Los tres son entornos esenciales para que la IA pase de generar lenguaje a hacer descubrimientos reales. El Agent de EP139, el laboratorio de modelos de EP138 y AI for Math de EP137 buscan, en el fondo, entornos que proporcionen retroalimentación.
\end{knowledgebox}

\subsection{El destino final de la colaboración entre humanos y máquinas}

El destino ideal de AI for Math no es necesariamente «sustituir a los matemáticos». Una forma más realista y también más valiosa es que la IA ayude a los matemáticos a formalizar demostraciones engorrosas, buscar lemma, encontrar contraejemplos, organizar la bibliografía y proponer conjeturas candidatas; los matemáticos, por su parte, se encargan del criterio para elegir problemas, la elección de direcciones, la interpretación y el juicio final. Así, la IA se convierte en una compañera de investigación y no en un simple sistema que responde preguntas de forma automática.

\begin{importantbox}{De herramienta a compañera}
Cuando la IA solo realiza pasos locales, es una herramienta; cuando puede comprender el contexto, acumular experiencia, proponer direcciones candidatas y someterse a verificación, empieza a parecerse a una compañera de investigación. AI for Math es un modelo de alta dificultad de esta transformación.
\end{importantbox}

\subsection{Resumen de la sección}

Las matemáticas son un campo de pruebas de alto valor para la entrada de la IA en el descubrimiento científico. Su fuerte verificabilidad hace que el entrenamiento sea más claro, y su estética e intuición nos recuerdan que el descubrimiento científico no puede reducirse a un único indicador.
\section{Mapa de lectura: de EP137 a los episodios siguientes de la serie}

EP137 ocupa un lugar particular en la serie de IA de Zhang Xiaojun. EP140 y EP138 se centran más en los laboratorios de modelos y en el paradigma industrial de los Agent; EP139 ofrece una historia técnica de los Agent; EP137, en cambio, lleva la cuestión hacia la investigación científica y la demostración formal. Es el puente entre «la IA transforma la productividad» y «la IA transforma la producción de conocimiento».

\subsection{Cinco preguntas que el lector debería llevarse}

Primera: ¿puede la IA pasar de resolver problemas a hacer verdadera investigación matemática? Segunda: ¿se convertirá Lean en la interfaz principal de AI for Math? Tercera: ¿puede un modelo aprender la intuición de los matemáticos, o esta tendrá que aportarla el ser humano a largo plazo? Cuarta: ¿cómo puede un Neo Lab como Axiom encontrar productos intermedios entre la investigación básica y la comercialización? Quinta: ¿el éxito de AI for Math impulsará, a su vez, el avance de los algoritmos y de la teoría de la propia IA?

\begin{knowledgebox}{Relación entre este episodio y las notas posteriores}
Si más adelante se elaboran notas de otras entrevistas sobre AI for Science, robótica, modelos del mundo o conducción autónoma, puede reutilizarse el marco de este episodio: primero buscar un entorno verificable, después examinar la retroalimentación de los expertos humanos y, por último, ver si el modelo puede formar un ciclo cerrado de largo plazo. Las matemáticas son solo el tipo de escenario de este marco en el que la verificación es más rigurosa.
\end{knowledgebox}

\subsection{Resumen de la sección}

El valor de EP137 está en que amplía los límites de la serie de IA: de los productos Agent y el entrenamiento de modelos pasa a las matemáticas y al descubrimiento científico. Nos recuerda que las aplicaciones de IA más difíciles suelen requerir, al mismo tiempo, verificación formal y sentido estético humano.
\section{Corrección de la transcripción y normalización de términos}

Este episodio no tiene subtítulos de la plataforma; el texto se basa en una transcripción hecha con faster-whisper large-v3 en CUDA. Los términos de matemáticas y de IA son abundantes, y la transcripción automática produce errores sistemáticos: por ejemplo, escribe Lean como «令» o «林», escribe AI for Math con palabras de sonido parecido, oye mal nombres propios como Carina, Hong o Verve Coffee y fragmenta términos como theorem proving, proof assistant y formalization. Al generar los apuntes es obligatorio corregir los términos; de lo contrario, el lector malinterpretará la línea técnica central.

\subsection{Términos clave: tabla de corrección de errores comunes de transcripción}

\noindent\begin{longtable}{p{0.24\textwidth}p{0.32\textwidth}p{0.35\textwidth}}
\toprule
Término normalizado & Posibles errores de transcripción & Motivo de la corrección \\
\midrule
Carina Hong / Hong Letong (洪乐潼) & 洪乐同, 洪乐童 y grafías erróneas de Carina por sonido parecido & El nombre está ligado a los antecedentes de Axiom; debe unificarse. \\
Axiom & 公理 («axioma»), A型 («tipo A»), Axium & Nombre de la empresa; en chino puede explicarse como «axioma», pero en el texto debe conservarse en inglés. \\
Lean & 令, 林, line & Nombre del proof assistant; está ligado a la línea central de la matemática formalizada. \\
AI for Math & Al for Math, AI for Mark & Dirección central de este episodio. \\
Proof Assistant & 证明助手 («asistente de demostración») y proof system usados indistintamente & Se refiere a sistemas de verificación como Lean o Coq. \\
Formalization & Confusión entre 形式化 («formalización») y 格式化 («dar formato») & Se refiere a escribir las matemáticas en un sistema formal, no a dar formato tipográfico. \\
Ken Ono & 小野肯, 小野 Ken & Nombre propio de un matemático, relacionado con las señales sobre la organización de Axiom. \\
Bounded / Free Attention & Traducciones erróneas como «límite de atención» o «atención libre» & Aquí es una metáfora de modos de pensamiento matemático, no el mecanismo de attention del Transformer en sí. \\
\bottomrule
\end{longtable}

\begin{warningbox}{La transcripción automática no puede usarse directamente como texto técnico}
El material de entrevistas es especialmente propenso a errores de transcripción en nombres propios, títulos de artículos y nombres de sistemas. Unas notas de alto estándar deben corregir los términos a partir de los capítulos, del contexto y de fuentes públicas externas; no pueden copiar palabra por palabra los subtítulos automáticos.
\end{warningbox}

\subsection{Resumen de la sección}

La transcripción es solo materia prima, no son los apuntes. Un contenido con tantos términos como AI for Math requiere especialmente normalizar primero los términos y después escribir notas estructuradas.
\section{La ruta de evaluación de AI for Math: del banco de problemas a la contribución a la investigación}

Para saber si AI for Math avanza de verdad, no basta con contar cuántos problemas de competencia resuelve el modelo. Los problemas de competencia tienen valor porque sus respuestas son claras y su retroalimentación es estricta, pero son solo el punto de entrada. Una evaluación de nivel superior debería considerar la tasa de aprobación de las demostraciones formales, las contribuciones a la biblioteca de teoremas, el descubrimiento de lemmas, la construcción de contraejemplos, la eficiencia de la colaboración con matemáticos y la aparición de nuevas líneas de investigación.

\subsection{Niveles de evaluación}

\noindent\begin{longtable}{p{0.22\textwidth}p{0.35\textwidth}p{0.34\textwidth}}
\toprule
Nivel & Indicadores medibles & Limitaciones \\
\midrule
Resolución de problemas de competencia & Tasa de aciertos, pass@k, tiempo de resolución & Es fácil sobreajustarse al banco de problemas; no refleja capacidad de investigación. \\
Completado formal & Tasa de aprobación de Lean theorem, tasa de éxito de tactic & Depende de la cobertura de la biblioteca de teoremas y de la selección de problemas. \\
Generación de demostraciones & Longitud de la demostración completa, legibilidad, tasa de aprobación de la verificación & Correcta no equivale a elegante; una demostración larga puede no aportar ninguna idea de fondo. \\
Contribución a la biblioteca de teoremas & Número de lemmas nuevos, tasa de reutilización, costo de mantenimiento & Requiere la aceptación de la comunidad; no basta con que la máquina los genere. \\
Colaboración en investigación & Tiempo que ahorran los matemáticos, contraejemplos encontrados, conjeturas propuestas & Es difícil de estandarizar, pero es lo más cercano al valor real. \\
\bottomrule
\end{longtable}

\begin{importantbox}{El núcleo de la ruta de evaluación}
La evaluación de AI for Math debería pasar gradualmente de «la respuesta es correcta» a «la demostración es verificable», «la demostración tiene poder explicativo», «la biblioteca de teoremas es reutilizable» y «los matemáticos están dispuestos a colaborar». Cuanto más se avanza, más difícil es cuantificar, pero también más cerca se está del valor real.
\end{importantbox}

\subsection{Por qué esto se relaciona con los Agents}

AI for Math parece un problema científico, pero en realidad también es un problema de Agents: el modelo tiene que entrar en el entorno de Lean, leer el goal state, invocar tactics, observar los errores, corregir la demostración y seguir avanzando. Este ciclo cerrado se parece mucho al de un web agent o un coding agent; la diferencia es que el entorno es más formal y la retroalimentación más estricta. Por eso, EP137 puede verse como una instancia científica de la historia técnica de los Agents que aborda EP139.

\begin{knowledgebox}{El entorno de Lean es un entorno de Agents de alta calidad}
Lean ofrece estados definidos, acciones definidas, verificación definida y retroalimentación de errores definida. Comparado con un entorno web abierto, es más riguroso; comparado con una conversación ordinaria, es más fácil de entrenar. Esto convierte a AI for Math en un campo de pruebas ideal para estudiar la capacidad de aprendizaje de los Agents.
\end{knowledgebox}

\subsection{Resumen de la sección}

La evaluación de AI for Math no puede quedarse en la tasa de aciertos sobre un banco de problemas. La ruta real avanza desde la resolución de problemas, la formalización y la generación de demostraciones hacia la contribución a la biblioteca de teoremas y la colaboración con matemáticos.

\section{Apéndice: ruta de aprendizaje de Lean y AI for Math}

\noindent\begin{longtable}{p{0.20\textwidth}p{0.35\textwidth}p{0.36\textwidth}}
\toprule
Etapa & Qué aprender & Por qué aprenderlo \\
\midrule
Lenguaje matemático & Proposiciones, definiciones, lemma, proof, contraejemplos & Entender cómo se organiza una demostración en lenguaje natural. \\
Fundamentos de Lean & theorem, example, rw, simp, apply, exact, induction & Entender cómo se interactúa con un proof assistant. \\
Mathlib & Biblioteca de teoremas de uso común, nomenclatura, búsqueda, relaciones de dependencia & Las demostraciones con IA dependen de las bibliotecas matemáticas existentes. \\
Autoformalization & Del lenguaje natural a un statement de Lean & Conectar artículos y enunciados de problemas con el sistema formal. \\
Proof Search & Combinación de tactics, goal state, corrección de errores & Permitir que el modelo itere con la retroalimentación del verificador. \\
Human Collaboration & Criterio estético de los matemáticos, elección de problemas, capacidad explicativa & Juzgar si una demostración tiene sentido, y no solo si pasa la verificación. \\
\bottomrule
\end{longtable}

\begin{knowledgebox}{El sentido de la ruta de aprendizaje}
Esta tabla cambia el orden de aprendizaje de AI for Math: en lugar de «primero aprender los modelos», propone «primero entender las matemáticas y los sistemas formales». Para los lectores con formación en ingeniería, los fundamentos de Lean y el proof search son el requisito mínimo para entrar en este campo; para los lectores con formación matemática, la autoformalization y la retroalimentación de los modelos son la vía de entrada para entender los límites de las capacidades de la IA.
\end{knowledgebox}

\subsection{Resumen de la sección}

AI for Math es una ruta interdisciplinaria: el lenguaje matemático, los sistemas formales, el entrenamiento de modelos y la colaboración entre humanos y máquinas son todos indispensables. Cuanto más clara sea la ruta de aprendizaje, menos fácil será malinterpretarla como «hacer que el modelo resuelva cada vez más ejercicios».
\section{Síntesis y ampliación}

\subsection{Conclusiones principales}

\begin{enumerate}
\item AI for Math no consiste en resolver ejercicios de matemáticas en serie, sino en un sistema completo que va del problema y la intuición a la formalización y la verificación.
\item Los proof assistants como Lean/Coq le dan a la IA una retroalimentación estricta, lo que convierte a las matemáticas en un dominio de alto valor, adecuado para el entrenamiento y la evaluación.
\item La intuición matemática y la verificación formal no se oponen; un sistema sólido necesita contar a la vez con free attention y bounded attention.
\item Un Neo Lab como Axiom solo se sostiene si confluyen una red de matemáticos, un sistema de formalización, el entrenamiento de modelos, financiamiento y capacidad de cómputo, y un objetivo de producto.
\item La meta de AI for Math no es sustituir a los matemáticos, sino construir junto con ellos demostraciones, bibliotecas de teoremas y líneas de investigación.
\end{enumerate}

\subsection{Preguntas abiertas}

\begin{itemize}
\item ¿Podrá la IA generar demostraciones con verdadero valor estético y poder explicativo, y no solo demostraciones largas y correctas?
\item ¿La formalización en Lean se convertirá en la interfaz principal de la investigación matemática, o solo servirá a algunas áreas?
\item ¿Cómo puede escalarse la retroalimentación de los matemáticos para que se forme un volante de datos de AI for Math?
\item ¿El producto de un Neo Lab como Axiom tomará la forma de un proof assistant, de una plataforma de investigación o de un sistema de colaboración entre matemáticos?
\end{itemize}

\subsection{Lecturas adicionales}

\begin{itemize}
\item Lean theorem prover y mathlib: para entender la infraestructura de las matemáticas formalizadas.
\item Las discusiones públicas de Terence Tao, Kevin Buzzard y otros sobre formalization: para entender por qué la comunidad matemática le da importancia a Lean.
\item Artículos sobre AI theorem proving, autoformalization y proof search: para entender las rutas técnicas de AI for Math.
\item Las notas de EP139 y EP138: para situar AI for Math dentro del panorama más amplio de los agentes y los laboratorios de modelos.
\end{itemize}

\begin{importantbox}{Valoración final}
La dificultad de AI for Math reside en que exige al modelo, al mismo tiempo, entender matemáticas, escribir en un lenguaje formal, someterse a la verificación automática y colaborar con la intuición de los matemáticos. Es el campo de prueba decisivo para que la IA pase de ser una «herramienta de productividad» a ser una «compañera en el descubrimiento científico».
\end{importantbox}

\end{document}
